\documentclass[12pt,a4paper]{article}
\usepackage[T2A]{fontenc}
\usepackage[utf8]{inputenc}
\usepackage[russian]{babel}
\usepackage{amsmath,amssymb}
\usepackage{geometry}
\geometry{margin=2.5cm}

\begin{document}

\begin{center}
{\Large\textbf{Алгоритм пространственно-морфологического анализа дендритных шипиков}}\\[0.5em]
{\normalsize Полное теоретическое и вычислительное описание конвейера анализа}
\end{center}

\noindent\textbf{Постановка задачи.}
Цель алгоритма состоит в количественном описании пространственной организации дендритных шипиков на поверхности дендрита и в последующем анализе связи этой организации с морфологическими свойствами шипиков. Для одного дендритного фрагмента рассматривается множество шипиков \(S=\{s_1,\ldots,s_n\}\). Каждому шипику сопоставляется точка крепления \(a_i\in\mathbb{R}^3\), точка геометрического центра \(c_i\in\mathbb{R}^3\), а также набор морфологических признаков \(m_i=(m_i^{(1)},\ldots,m_i^{(p)})\). В текущей массовой версии анализа основными признаками шипика являются объём \(V_i\) и длина \(L_i\), поскольку они устойчиво интерпретируются биологически и используются в последующих процедурах сравнения, автокорреляции и кластерного анализа.

\noindent\textbf{Загрузка данных и формирование объекта дендрита.}
Вычислительный конвейер начинается с загрузки датасета и создания экземпляра класса \texttt{Dendrite}. Для формата \texttt{labid} используется исходная логика чтения папки с мешем дендрита и соответствующими мешами шипиков. Для формата \texttt{microns} анализируются папки нейронов, внутри которых выбираются ветви вида \texttt{limb\_\{id\}/branch\_\{id\}}; целевым мешем дендрита является \texttt{branch\_mesh.off}, а целевые меши шипиков считываются из папки \texttt{spines}. Ветка исключается из анализа, если папка со шипиками отсутствует или если не удаётся корректно загрузить необходимые меши. После чтения данных объект \texttt{Dendrite} хранит геометрию дендрита, набор мешей шипиков, координаты точек крепления, координаты центров шипиков и вычисленные метрики.

\noindent\textbf{Первичные геометрические характеристики дендрита.}
При инициализации дендрита вычисляются базовые характеристики самого дендритного фрагмента: объём \(V_D\), длина \(L_D\) и эффективный радиус \(R_D\). Объём по возможности вычисляется по мешу; если прямой расчёт недоступен или меш не подходит для CGAL-процедур, используется fallback на основе \texttt{trimesh}. Длина дендрита предпочтительно берётся из зарегистрированного скелета, если он доступен. Если скелет не зарегистрирован, выполняется скелетизация для \texttt{Polyhedron\_3}; если она невозможна, используется приближённая оценка длины по центральной оси меша. Эффективный радиус вычисляется из цилиндрической аппроксимации
\[
R_D=\sqrt{\frac{V_D}{\pi L_D}}.
\]
Эта величина не утверждает, что дендрит является идеальным цилиндром, но даёт грубый геометрический масштаб фрагмента. Внутри класса также могут вычисляться служебные величины \(dr\) и \(V_{\mathrm{around}}\), где \(dr\) равен половине максимальной длины шипика, а \(V_{\mathrm{around}}=\pi L_D((R_D+dr)^2-R_D^2)\) соответствует объёму цилиндрической оболочки вокруг дендрита. Эти величины использовались в старой трёхмерной нормировке \texttt{NNdist\_norm}, но не включаются в актуальный итоговый вектор признаков, поскольку анализ шипиков сейчас основан на расстояниях по поверхности дендрита, а не на объёмной 3D-плотности в оболочке.

\noindent\textbf{Валидация и морфологические метрики шипиков.}
Перед включением шипика в анализ выполняется проверка корректности его меша. Меш должен иметь конечные координаты вершин, набор граней и хотя бы одну невырожденную грань. Дополнительно применяется эвристика отбраковки объектов, которые вероятно являются не шипиками, а ошибочно выделенными фрагментами ствола дендрита. Такой фрагмент часто имеет форму трубки с двумя открытыми концами сопоставимого размера; поэтому анализируются граничные компоненты меша, их радиусы, относительные размеры и вытянутость объекта вдоль главной оси. Если объект похож на трубчатый фрагмент дендрита, он не используется как шипик. Для валидных шипиков вычисляются морфологические признаки; в текущей облегчённой конфигурации сохраняются объём и длина, тогда как остальные морфологические метрики временно исключены из массового расчёта для уменьшения вычислительной нагрузки.

\noindent\textbf{Определение точек крепления.}
Для каждого шипика определяется точка крепления \(a_i\), которая используется как пространственная позиция шипика на поверхности дендрита. Если доступен внешний файл с заранее вычисленными точками крепления, используется он. Если такой файл отсутствует, точка определяется по геометрии меша шипика. В случае \texttt{trimesh}-мешей находятся граничные рёбра, то есть рёбра, принадлежащие ровно одной грани. Эти рёбра разбиваются на связные компоненты, каждая из которых рассматривается как потенциальное граничное кольцо основания шипика. Компонента считается замкнутой, если каждая её вершина имеет степень два в графе граничных рёбер. Если замкнутых компонент несколько, выбирается крупнейшая по числу рёбер; если замкнутых компонент нет, используется крупнейшая граничная компонента как fallback. Точка крепления вычисляется как центр масс вершин выбранного кольца:
\[
a_i=\frac{1}{|B_i|}\sum_{v\in B_i}x_v,
\]
где \(B_i\) обозначает множество вершин выбранной граничной компоненты, а \(x_v\) есть координаты вершины. Такая процедура соответствует геометрическому предположению, что основание шипика после сегментации образует наиболее выраженную открытую границу.

\noindent\textbf{Переход к расстояниям по поверхности дендрита.}
Центральным объектом пространственного анализа является матрица расстояний \(D=(d_{ij})\), где \(d_{ij}\) есть кратчайшее расстояние между точками крепления \(a_i\) и \(a_j\) по сетке дендрита. Меш дендрита рассматривается как взвешенный граф \(G_M=(V_M,E_M)\), вершины которого соответствуют вершинам треугольной сетки, а рёбра соответствуют сторонам треугольников. Вес ребра равен евклидовой длине соответствующего сегмента. Каждая точка крепления проецируется на ближайшую вершину сетки, после чего расстояние между двумя шипиками определяется как длина кратчайшего пути:
\[
d_{ij}=\min_{\pi:i\rightarrow j}\sum_{(u,v)\in\pi}\|x_u-x_v\|_2.
\]
Использование \texttt{mesh\_graph}-расстояния принципиально важно, поскольку оно измеряет близость вдоль поверхности дендрита, а не прямолинейную близость в трёхмерном пространстве. Для изогнутых ветвей это предотвращает ошибочную интерпретацию двух геометрически близких, но топологически далёких участков как соседних.

\noindent\textbf{Этап оценки плотности и общего распределения шипиков.}
Первый аналитический этап после построения матрицы расстояний реализован функцией \texttt{calculate\_density\_distribution\_analysis}. На этом этапе вычисляются метрики плотности и общего распределения шипиков без предположения о наличии кластеров. Для каждого шипика определяется расстояние до ближайшего соседа, а также расстояния до второго, третьего и пятого ближайших соседей. Если \(d_{i(1)}\leq d_{i(2)}\leq\ldots\) обозначает упорядоченные положительные расстояния от шипика \(s_i\) до остальных шипиков, то \(d_{i(1)}\) характеризует локальную разреженность около данного шипика. Для всего дендрита сохраняются среднее, медиана, стандартное отклонение, коэффициент вариации и квантили 0.10, 0.25, 0.75 и 0.90 nearest-neighbor расстояний:
\[
\mathrm{CV}_{NN}=\frac{\mathrm{sd}(d_{i(1)})}{\mathrm{mean}(d_{i(1)})}.
\]
Низкие значения nearest-neighbor расстояний соответствуют плотному расположению шипиков, а высокий коэффициент вариации указывает на неоднородное распределение, при котором плотные участки сосуществуют с разреженными.

\noindent\textbf{Профиль попарных расстояний и энтропия пространственного распределения.}
В рамках того же этапа вычисляется профиль попарных расстояний, или \texttt{pair\_distance\_profile}. Для набора радиусов \(r\) и ширины бина \(dr\) оценивается среднее число соседей на один шипик, расстояния до которых попадают в интервал \([r,r+dr)\). В дискретной форме это записывается как
\[
h(r)=\frac{1}{n}\sum_{i\neq j}\mathbf{1}\{r\leq d_{ij}<r+dr\}.
\]
Эта величина не является нормированной pair correlation function в строгом статистическом смысле, поскольку она не делится на ожидаемое число пар при случайном расположении шипиков на том же дендрите. Поэтому она интерпретируется как описательная характеристика распределения попарных расстояний. По значениям \(h(r)\) вычисляется энтропия Шеннона:
\[
H=-\sum_k p_k\log p_k,
\]
где \(p_k\) есть нормированное значение профиля в \(k\)-м интервале расстояний. Если сумма значений профиля равна нулю, энтропия считается равной нулю. Такая энтропия служит интегральной мерой разнообразия пространственных масштабов: более равномерное распределение пар по радиальным интервалам даёт большую энтропию, а концентрация пар в узком диапазоне расстояний уменьшает её. В актуальный итоговый вектор дендрита включается только энтропия профиля, тогда как сами массивы значений профиля и радиусов не сохраняются как признаки для обучения модели. Настоящая PCF может быть добавлена как отдельная метрика, если для дендрита будет задана null-модель случайного расположения шипиков.

\noindent\textbf{Попарные расстояния.}
В рамках этапа общего распределения сохраняется распределение всех попарных расстояний \(d_{ij}\), \(i<j\). Для него вычисляются среднее, медиана, стандартное отклонение, коэффициент вариации и квантили 0.10, 0.25, 0.75 и 0.90. Коэффициент вариации
\[
\mathrm{CV}_{\mathrm{pair}}=\frac{\mathrm{sd}(d_{ij}:i<j)}{\mathrm{mean}(d_{ij}:i<j)}
\]
отражает относительную неоднородность общего распределения расстояний между всеми парами шипиков: высокие значения соответствуют ситуации, когда на одном фрагменте присутствуют как близкие пары, так и пары, разделённые большим расстоянием.

\noindent\textbf{Выбор локального масштаба соседства.}
Для дальнейшего локального анализа требуется радиус соседства \(\varepsilon\). Если DBSCAN уже был запущен и имеет положительный радиус, используется его \(\varepsilon\). Если радиус отсутствует или некорректен, он выбирается по k-distance эвристике: \(\texttt{min\_samples}=3\), а \(\varepsilon\) равен 75-му перцентилю расстояний до третьего ближайшего соседа:
\[
\varepsilon=Q_{0.75}\bigl(\{d_{i(3)}\}_{i=1}^{n}\bigr).
\]
Практически это означает, что локальный масштаб адаптируется к конкретной ветви дендрита и задаёт такой радиус, внутри которого у типичного шипика в достаточно плотной области должны находиться несколько соседей. Этот радиус используется не только для DBSCAN, но и для построения порогового графа соседства.

\noindent\textbf{Этап анализа локальных окрестностей шипиков.}
Второй аналитический этап реализован функцией \texttt{calculate\_local\_neighborhood\_analysis}. На выбранном масштабе \(\varepsilon\) строится граф \(G_\varepsilon\), вершины которого соответствуют шипикам, а ребро \((i,j)\) проводится тогда и только тогда, когда \(d_{ij}\leq\varepsilon\). Окрестность шипика \(s_i\) определяется как множество всех шипиков, лежащих не дальше \(\varepsilon\):
\[
N_\varepsilon(i)=\{j:j\neq i,\ d_{ij}\leq\varepsilon\}.
\]
Такие окрестности могут пересекаться, поскольку один и тот же шипик может входить в окрестности нескольких центральных шипиков. Локальное число соседей шипика равно
\[
\rho_i=\sum_{j\neq i}\mathbf{1}\{d_{ij}\leq\varepsilon\}.
\]
Для каждого шипика сохраняется \(\rho_i\), степень в пороговом графе, локальный коэффициент кластеризации, средние морфологические значения соседей и абсолютное отличие собственной морфологии от соседнего среднего. Для всего дендрита сохраняются среднее и максимальное число соседей, доля шипиков без соседей, число рёбер в \(G_\varepsilon\), число компонент связности и размер крупнейшей компоненты. На этом же масштабе вычисляется доля близких пар
\[
F_\varepsilon=\frac{\#\{(i,j):i<j,\ d_{ij}\leq\varepsilon\}}{\#\{(i,j):i<j\}}.
\]
Эта величина является простой Ripley-like характеристикой на фиксированном масштабе: чем она выше, тем большая часть всех пар шипиков располагается в пределах локального радиуса соседства \(\varepsilon\). В отличие от полной функции Рипли, такая оценка не требует моделирования случайных точек на сложной поверхности дендрита и поэтому удобна для массовой обработки. Этот этап является актуальной версией старого анализа пересекающихся окрестностей: окрестность каждого шипика по-прежнему рассматривается как локальная группа, но результаты сохраняются как признаки вершин и сводные характеристики порогового графа, а не как отдельный набор пересекающихся DBSCAN-подобных кластеров.

\noindent\textbf{Этап морфологической автокорреляции.}
Третий аналитический этап реализован функцией \texttt{calculate\_spatial\_autocorrelation\_analysis}. Внутри этого этапа вычисляются пространственные статистики для всех доступных числовых морфологических признаков, включая объём и длину. Для графа соседства с весами \(w_{ij}\), где \(w_{ij}=1\) при \(d_{ij}\leq\varepsilon\) и \(w_{ij}=0\) иначе, глобальный индекс Морана \(I\) определяется как
\[
I=\frac{n}{W}\frac{\sum_{i,j}w_{ij}(x_i-\bar x)(x_j-\bar x)}{\sum_i(x_i-\bar x)^2},
\qquad W=\sum_{i,j}w_{ij}.
\]
Положительные значения \(I\) указывают на пространственное соседство похожих морфологических значений, отрицательные значения на соседство контрастных значений, а значения около нуля соответствуют отсутствию выраженной пространственной структуры.

\noindent\textbf{Коэффициент Гири \(C\), соседская корреляция Спирмена и перестановочная проверка.}
Для дополнительной проверки морфологической структуры вычисляется коэффициент Гири \(C\):
\[
C=\frac{n-1}{2W}\frac{\sum_{i,j}w_{ij}(x_i-x_j)^2}{\sum_i(x_i-\bar x)^2}.
\]
Значения \(C<1\) соответствуют положительной автокорреляции, \(C\approx1\) отсутствию пространственной структуры, а \(C>1\) соседству непохожих значений. Значения соседних средних, вычисленные на этапе локальных окрестностей, используются для дополнительной соседской корреляции Спирмена:
\[
\bar m_{N(i)}^{(r)}=\frac{1}{|N_\varepsilon(i)|}\sum_{j:d_{ij}\leq\varepsilon}m_j^{(r)}
\]
После этого оценивается ранговая корреляция Спирмена между собственным значением признака и средним значением признака у соседей. Значимость глобального индекса Морана \(I\), коэффициента Гири \(C\) и соседней корреляции Спирмена проверяется перестановочно: значения морфологического признака случайно переставляются между фиксированными пространственными позициями, что разрушает связь между морфологией и расположением, но сохраняет само распределение признака. Перестановочное \(p\)-значение вычисляется как
\[
p=\frac{1+\#\{|\theta_b|\geq|\theta_{\mathrm{obs}}|\}}{B+1},
\]
где \(\theta_{\mathrm{obs}}\) есть наблюдаемая статистика, \(\theta_b\) статистика на \(b\)-й перестановке, а \(B\) число перестановок.

\noindent\textbf{Этап пространственной кластеризации DBSCAN.}
Четвёртый аналитический этап начинается с функции \texttt{calculate\_cluster\_metrics}. Кластеризация выполняется алгоритмом DBSCAN с предвычисленной матрицей \texttt{mesh\_graph}-расстояний. Параметр \(\texttt{min\_samples}\) фиксирован равным трём, а радиус \(\varepsilon\) выбирается по 75-му перцентилю расстояний до третьего ближайшего соседа, если он не задан явно. DBSCAN присваивает каждому шипику метку кластера \(c_i\) или метку шума \(-1\). Доля шумовых объектов
\[
q_{\mathrm{noise}}=\frac{\#\{i:c_i=-1\}}{n}
\]
интерпретируется как доля изолированных или пространственно разреженных шипиков. В отличие от простого порогового графа, DBSCAN выделяет плотностно-связанные области и отделяет точки, которые не принадлежат ни одной плотной группе.

\noindent\textbf{Сводные характеристики пространственных кластеров.}
Функция \texttt{calculate\_spatial\_cluster\_analysis} описывает найденные DBSCAN-кластеры. Для каждого кластера \(C_k=\{i:c_i=k\}\) вычисляется размер, диаметр
\[
\Delta_k=\max_{i,j\in C_k}d_{ij},
\]
среднее и медианное внутрикластерное расстояние, а также линейная плотность \(|C_k|/\Delta_k\), если \(\Delta_k>0\). Для каждой морфологической метрики внутри кластера вычисляются среднее, медиана и стандартное отклонение. Эти величины позволяют понять, являются ли пространственно плотные группы морфологически однородными и отличаются ли отдельные кластеры по объёму или длине шипиков.

\noindent\textbf{Сравнение кластеризованных и изолированных шипиков.}
После DBSCAN шипики разделяются на кластеризованные, для которых \(c_i\neq -1\), и изолированные, для которых \(c_i=-1\). Для каждой морфологической метрики сравниваются распределения этих двух групп. В качестве описательных характеристик для каждой группы сохраняются среднее, медиана, стандартное отклонение и коэффициент вариации. Среднее отражает общий уровень признака в группе, медиана задаёт устойчивую к выбросам характеристику типичного значения, стандартное отклонение описывает абсолютный разброс, а коэффициент вариации нормирует этот разброс на средний уровень признака. Для проверки различий используются тест Манна--Уитни и тест Бруннера--Мунцеля, поскольку распределения морфологических признаков могут быть ненормальными и асимметричными. Кроме \(p\)-значений вычисляется размер эффекта Cliff's delta:
\[
\delta=\frac{\#\{(x,y):x>y\}-\#\{(x,y):x<y\}}{n_xn_y},
\]
где \(x\) пробегает значения признака среди кластеризованных шипиков, а \(y\) среди изолированных. Значение \(\delta>0\) означает, что признак чаще выше у кластеризованных шипиков, \(\delta<0\) означает, что он чаще выше у изолированных, а значение около нуля указывает на слабый эффект.

\noindent\textbf{Графовый анализ пространственной организации.}
После кластерного и пространственно-морфологического анализа может выполняться функция \texttt{graph\_analysis}. Здесь строится взвешенный граф всех шипиков, где ребро между парой шипиков получает вес, обратно пропорциональный расстоянию \(1/d_{ij}\). На этом графе вычисляются средний коэффициент кластеризации, размеры сообществ, средний размер сообщества, характерная протяжённость сообщества и модулярность. Если доступен пакет \texttt{python-louvain}, сообщества находятся алгоритмом Louvain; иначе используется жадная оптимизация модулярности из \texttt{networkx}. Графовые метрики дополняют DBSCAN: DBSCAN выделяет локальные плотностные кластеры, а графовый анализ описывает более общую модульную организацию сети пространственных близостей.

\noindent\textbf{Общая организация кода после реорганизации.}
Алгоритм в коде разделён на логические этапы без дублирования публичных функций с близким смыслом. Функция \texttt{calculate\_density\_distribution\_analysis} отвечает за этап плотности и общего распределения: внутри неё вычисляются \texttt{NNdist}, \texttt{pair\_distance\_profile}, энтропия профиля попарных расстояний, nearest-neighbor статистики и pair-distance статистики. Функция \texttt{calculate\_local\_neighborhood\_analysis} отвечает за этап локальных окрестностей и порогового графа соседства. Функция \texttt{calculate\_spatial\_autocorrelation\_analysis} отвечает за этап пространственной морфологической автокорреляции: внутри неё вычисляются глобальный индекс Морана \(I\), коэффициент Гири \(C\), соседская корреляция Спирмена и перестановочные \(p\)-значения для каждой доступной числовой морфологической метрики. Функция \texttt{calculate\_spatial\_cluster\_analysis} отвечает за описание DBSCAN-кластеров и сравнение кластеризованных и изолированных шипиков. Функция \texttt{calculate\_comprehensive\_spatial\_analysis} является оркестратором этих этапов и запускает их в порядке: плотность и распределение, локальные окрестности, автокорреляция, кластерный анализ.

\noindent\textbf{Сохраняемые результаты.}
Базовые характеристики дендрита сохраняются в \texttt{dendr\_metrics.json}, координаты шипиков в \texttt{spine\_coords.json}, морфологические метрики шипиков в \texttt{spine\_metrics.json}. Результаты этапа плотности сохраняются в \texttt{grouping\_dendr\_metrics.json}; там находятся \texttt{NNdist}, \texttt{NNdist\_norm}, значения \texttt{PairDistanceProfile}, радиусы профиля и энтропия профиля, что сохраняет возможность просматривать старые диагностические файлы. Результаты DBSCAN сохраняются в \texttt{cluster\_dendr\_metrics.json}, графовые метрики в \texttt{graph\_dendr\_metrics.json}. Расширенный анализ сохраняет таблицы \texttt{spatial\_morphology\_summary.csv}, \texttt{spatial\_morphology\_spines.csv}, \texttt{spatial\_morphology\_clusters.csv}, \texttt{spatial\_morphology\_tests.csv} и \texttt{spatial\_morphology\_permutation.csv}. Первая таблица содержит одну строку на дендрит и описывает глобальную плотность, nearest-neighbor статистики, долю шума, долю близких пар и характеристики порогового графа. Вторая содержит одну строку на шипик и описывает его локальное окружение, DBSCAN-метку, расстояния до соседей и морфологическое отличие от соседей. Третья содержит одну строку на DBSCAN-кластер, четвёртая содержит результаты сравнения кластеризованных и изолированных шипиков, а пятая содержит статистики пространственной морфологической автокорреляции и перестановочные \(p\)-значения. Кроме детальных файлов формируется общий файл \texttt{dendrite\_structural\_organization\_vectors.csv}, в котором одна строка соответствует одному дендритному фрагменту, а столбцы образуют единый вектор структурной организации дендрита. В этот вектор включаются только агрегированные скалярные признаки, пригодные для последующего анализа и обучения модели; подробные переменные по отдельным кластерам, массивы профиля попарных расстояний и размеры отдельных графовых сообществ остаются вне общего ML-вектора.

\noindent\textbf{Интерпретация полного конвейера.}
Полный алгоритм последовательно переходит от геометрии мешей к биологически интерпретируемым пространственным характеристикам. Сначала определяется корректный набор шипиков и их точки крепления к дендриту. Затем строится матрица расстояний по поверхности дендрита, которая становится общей основой для всех последующих расчётов. Далее оценивается общая плотность расположения шипиков через nearest-neighbor статистики, профиль попарных расстояний, энтропию этого профиля и pair-distance распределение. После этого анализируются локальные пересекающиеся окрестности шипиков и пороговый граф соседства. Затем проверяется, связана ли морфология шипиков с их пространственным окружением, с помощью глобального индекса Морана \(I\), коэффициента Гири \(C\), соседской корреляции Спирмена и перестановочных проверок. Далее выделяются плотностные группы шипиков алгоритмом DBSCAN, описывается геометрия этих групп и проверяется, отличаются ли кластеризованные шипики от изолированных по объёму и длине. Наконец, графовые метрики дают дополнительное описание модульной структуры пространственных связей. Таким образом, результатом является не только набор отдельных метрик, а многоуровневое описание взаимного расположения шипиков, плотности их распределения, локальной кластеризации и связи пространственной организации с морфологией.

\end{document}
